% Vietnamese translation for OI Public Library
% Source-PDF: 比赛 CONTEST/2021“MINIEYE杯”中国大学生算法设计超级联赛/2021“MINIEYE杯”中国大学生算法设计超级联赛（6）/2021“MINIEYE杯”中国大学生算法设计超级联赛（6）-题目集.pdf
\documentclass[11pt]{article}
\usepackage{oipl-vn}

\title{Đề bài trận thứ sáu HDU Multi-University 2021}
\author{}
\date{}

\begin{document}
\maketitle

\origtitle{2021 HDU Multi-University Training Contest 6 (XJTU Contest)}
\sourcepath{contest/hdu-2021-minieye-06-problems.pdf}

\section{Problem A. Yes, Prime Minister}

\paragraph{Tệp vào:} standard input
\paragraph{Tệp ra:} standard output
\paragraph{Giới hạn thời gian:} 6 seconds
\paragraph{Giới hạn bộ nhớ:} 512 megabytes

Bộ Các vấn đề Hành chính của ngài Hacker có vô hạn công chức. Mỗi số nguyên
được dùng làm mã định danh bởi đúng một công chức. Ngài Hacker rất muốn giảm
biên chế trong bộ máy công vụ, nên ông sẽ chỉ giữ lại những người có mã định
danh liên tiếp trong đoạn \([l,r]\) và sa thải những người còn lại.

Tuy nhiên, mã định danh của thư ký thường trực Sir Humphrey là \(x\), và ông
không thể bị đuổi việc, nên phải có \(l\le x\le r\). Ngài Hacker muốn trở
thành Thủ tướng, vì vậy ông yêu cầu tổng mã định danh của những người được
giữ lại,
\[
  \sum_{i=l}^{r} i,
\]
phải là một số nguyên tố.

Bạn, Bernard, cần lập phương án cắt giảm thỏa mãn yêu cầu của cả hai cấp
trên. Nếu không, ngài Hacker hoặc Sir Humphrey sẽ sa thải bạn.

Ngài Hacker sẽ hài lòng nếu giữ lại càng ít người càng tốt. Hãy tính số
lượng người tối thiểu còn lại để đáp ứng các yêu cầu của họ.

Một số nguyên tố \(p\) là một số nguyên lớn hơn \(1\), không có ước nguyên
dương nào ngoài \(1\) và \(p\).

\subsection*{Dữ liệu vào}

Dòng đầu tiên chứa một số nguyên \(T\) (\(1\le T\le 10^6\)) là số bộ dữ
liệu. Sau đó là \(T\) bộ dữ liệu.

Dòng duy nhất của mỗi bộ dữ liệu chứa một số nguyên \(x_i\)
(\(-10^7\le x_i\le 10^7\)) là mã định danh của Sir Humphrey.

\subsection*{Dữ liệu ra}

Với mỗi bộ dữ liệu, hãy in số lượng người tối thiểu cần giữ lại nếu tồn tại
một phương án như vậy; nếu không, in \(-1\).

\subsection*{Ví dụ}

\paragraph{standard input}
\begin{verbatim}
10
-2
-1
0
1
2
3
4
5
6
7
\end{verbatim}

\paragraph{standard output}
\begin{verbatim}
6
4
3
2
1
1
2
1
2
1
\end{verbatim}

\section{Problem B. Might and Magic}

\paragraph{Tệp vào:} standard input
\paragraph{Tệp ra:} standard output
\paragraph{Giới hạn thời gian:} 12 seconds
\paragraph{Giới hạn bộ nhớ:} 512 megabytes

Hai anh hùng đang giao chiến, lần lượt có tên là anh hùng \(0\) và anh hùng
\(1\).

Bạn điều khiển anh hùng \(0\), còn kẻ địch của bạn là anh hùng \(1\). Mỗi
anh hùng có năm thuộc tính nguyên: ATTACK, DEFENCE, POWER, KNOWLEDGE và
HEALTH. Khi hai anh hùng giao đấu, họ sẽ luân phiên tấn công và anh hùng của
bạn ra đòn trước. Trong một lượt, một anh hùng có thể thực hiện đúng một đòn
tấn công, hoặc tấn công vật lý hoặc tấn công phép thuật.

Giả sử các thuộc tính của họ là \(A_i,D_i,P_i,K_i,H_i\) (\(0\le i\le 1\)).
Với anh hùng \(i\), sát thương của đòn tấn công vật lý là
\(C_p\max(1,A_i-D_{1-i})\), còn sát thương của đòn tấn công phép thuật là
\(C_mP_i\), trong đó \(C_p,C_m\) là các hằng số cho trước. Anh hùng \(i\) có
thể tấn công phép thuật không quá \(K_i\) lần trong toàn bộ trận đấu. Sau khi
anh hùng \(i\) tấn công, \(H_{1-i}\) sẽ giảm đi lượng sát thương do đối thủ
gây ra. Nếu \(H_{1-i}\le 0\), anh hùng \(1-i\) thua và trận đấu kết thúc.

Bây giờ bạn biết kẻ địch của mình là Yog, một người hoàn toàn không biết
phép thuật, nghĩa là \(P_1=K_1=0\) và hắn sẽ chỉ tấn công vật lý. Bạn có thể
phân phối tùy ý \(N\) điểm thuộc tính vào bốn thuộc tính
\(A_0,D_0,P_0,K_0\), nghĩa là các thuộc tính này có thể là các số nguyên
không âm bất kỳ thỏa mãn \(0\le A_0+D_0+P_0+K_0\le N\).

Cho \(C_p,C_m,H_0,A_1,D_1\) và \(N\), hãy tính giá trị lớn nhất của \(H_1\)
mà bạn vẫn có cơ hội thắng.

\subsection*{Dữ liệu vào}

Dòng đầu tiên chứa một số nguyên \(T\) (\(1\le T\le 10^5\)) là số bộ dữ
liệu. Sau đó là \(T\) bộ dữ liệu.

Dòng duy nhất của mỗi bộ dữ liệu chứa sáu số nguyên
\(C_p,C_m,H_0,A_1,D_1,N\)
\[
  1\le C_p,C_m,H_0,A_1,D_1,N\le 10^6,
\]
là các thuộc tính đã mô tả ở trên.

\subsection*{Dữ liệu ra}

Với mỗi bộ dữ liệu, in một số nguyên biểu thị đáp án.

\subsection*{Ví dụ}

\paragraph{standard input}
\begin{verbatim}
2
1 1 4 5 1 4
2 5 1 9 9 6
\end{verbatim}

\paragraph{standard output}
\begin{verbatim}
4
25
\end{verbatim}

\section{Problem C. 0 tree}

\paragraph{Tệp vào:} standard input
\paragraph{Tệp ra:} standard output
\paragraph{Giới hạn thời gian:} 2 seconds
\paragraph{Giới hạn bộ nhớ:} 512 megabytes

Ta có một cây \(\langle V,E\rangle\) gồm \(n\) đỉnh, với trọng số \(a_i\)
(\(i\in V\)) trên các đỉnh và trọng số \(b_e\) (\(e=\langle u,v\rangle\in E\))
trên các cạnh hai chiều. \(a_i\) là số nguyên không âm, còn \(b_e\) là số
nguyên.

Bạn có thể thực hiện không quá \(4n\) thao tác sau.

Với đường đi ngắn nhất từ \(x\) đến \(y\) đi qua \(k+1\) đỉnh
\((v_0,v_1,v_2,\ldots,v_k)\), trong đó \(v_0=x\), \(v_k=y\), đặt
\(e_i=\langle v_i,v_{i+1}\rangle\) (\(0\le i<k\)). Trong một thao tác, bạn có
thể chọn hai đỉnh \(x,y\) và một số nguyên không âm \(w\), để thực hiện:
\[
  a_x\leftarrow a_x\oplus w,\qquad
  a_y\leftarrow a_y\oplus w,\qquad
  b_{e_i}\leftarrow b_{e_i}+(-1)^i w\quad (0\le i<k).
\]
Ở đây \(\oplus\) là phép XOR bit. Có thể thấy rằng nếu \(x=y\), sẽ không có
gì thay đổi.

Bạn cần quyết định liệu có thể biến mọi \(a_i,b_e\) thành \(0\) hay không.
Nếu có thể, hãy in một nghiệm dùng không quá \(4n\) thao tác như mô tả ở
trên.

\subsection*{Dữ liệu vào}

Dòng đầu tiên chứa một số nguyên \(T\) (\(1\le T\le 250\)) là số bộ dữ liệu.
Sau đó là \(T\) bộ dữ liệu.

Dòng đầu tiên của mỗi bộ dữ liệu chứa một số nguyên \(n\) (\(1\le n\le
10^4\)) là số đỉnh.

Dòng tiếp theo chứa \(n\) số nguyên không âm \(a_i\) (\(0\le a_i<2^{30}\)) là
trọng số trên mỗi đỉnh.

Sau đó là \(n-1\) dòng, mỗi dòng chứa ba số nguyên \(x_j,y_j,w_j\)
\[
  1\le x_j,y_j\le n,\qquad -10^9\le w_j\le 10^9,
\]
biểu thị một cạnh giữa \(x_j\) và \(y_j\) có trọng số \(w_j\). Bảo đảm rằng
các cạnh đã cho tạo thành một cây.

Bảo đảm rằng \(\sum n\le 10^5\).

\subsection*{Dữ liệu ra}

Với mỗi bộ dữ liệu, in \texttt{YES} ở dòng đầu tiên nếu bạn có thể biến mọi
\(a_i,b_e\) thành \(0\) bằng không quá \(4n\) thao tác. Nếu không, in
\texttt{NO}.

Nếu có thể biến mọi trọng số thành \(0\), hãy in nghiệm của bạn trong \(k+1\)
dòng tiếp theo (\(0\le k\le 4n\)) như sau.

Dòng đầu tiên chứa một số nguyên \(k\) là số thao tác bạn thực hiện.

Sau đó là \(k\) dòng, mỗi dòng chứa ba số nguyên \(X_j,Y_j,W_j\)
\[
  1\le X_j,Y_j\le n,\qquad 0\le W_j\le 10^{14},
\]
biểu thị một thao tác.

Nếu có nhiều nghiệm, in bất kỳ nghiệm nào.

\subsection*{Ví dụ}

\paragraph{standard input}
\begin{verbatim}
3
1
0
2
2 3
1 2 -2
3
5 6 3
1 2 -5
2 3 -5
\end{verbatim}

\paragraph{standard output}
\begin{verbatim}
YES
0
NO
YES
3
1 3 5
2 3 7
2 3 3
\end{verbatim}

\section{Problem D. Decomposition}

\paragraph{Tệp vào:} standard input
\paragraph{Tệp ra:} standard output
\paragraph{Giới hạn thời gian:} 2 seconds
\paragraph{Giới hạn bộ nhớ:} 512 megabytes

Cho một đồ thị đầy đủ vô hướng có \(n\) đỉnh, trong đó \(n\) là số lẻ. Bạn
cần phân hoạch tập cạnh của nó thành \(k\) đường đi đơn rời nhau, sao cho
đường đi đơn thứ \(i\) có độ dài \(l_i\) (\(1\le i\le k,\ 1\le l_i\le n-3\)),
và mỗi cạnh vô hướng được dùng đúng một lần.

Đồ thị đầy đủ là một đồ thị đơn vô hướng trong đó mọi cặp đỉnh phân biệt đều
được nối bởi đúng một cạnh. Một đường đi đơn có độ dài \(l\) ở đây nghĩa là
đường đi đi qua \(l\) cạnh, và các đỉnh trên đường đi đôi một phân biệt.

Có thể chứng minh rằng luôn tồn tại đáp án nếu
\[
  \sum_{i=1}^{k} l_i=\frac{n(n-1)}{2}.
\]

\subsection*{Dữ liệu vào}

Dòng đầu tiên chứa một số nguyên \(T\) (\(1\le T\le 100000\)) là số bộ dữ
liệu. Sau đó là \(T\) bộ dữ liệu.

Dòng đầu tiên của mỗi bộ dữ liệu chứa hai số nguyên \(n,k\)
\[
  5\le n\le 1000,\qquad
  1\le k\le \frac{n(n-1)}{2},\qquad
  n\equiv 1\pmod 2,
\]
là số đỉnh và số đường đi.

Dòng tiếp theo chứa \(k\) số nguyên \(l_1,l_2,\ldots,l_k\)
(\(1\le l_i\le n-3\)) là độ dài của từng đường đi.

Bảo đảm rằng \(\sum_{i=1}^{k}l_i=\frac{n(n-1)}{2}\) với mỗi bộ dữ liệu, và
\(\sum \frac{n(n-1)}{2}\le 4\times 10^6\).

\subsection*{Dữ liệu ra}

Với mỗi bộ dữ liệu, trước tiên in một dòng chứa \texttt{Case \#x:}, trong đó
\(x\) (\(1\le x\le T\)) là số thứ tự bộ dữ liệu. Sau đó in \(k\) dòng. Dòng
thứ \(i\) chứa \(l_i+1\) số biểu thị đường đi thứ \(i\).

Nếu có nhiều đáp án, in bất kỳ đáp án nào.

\subsection*{Ví dụ}

\paragraph{standard input}
\begin{verbatim}
3
5 6
2 1 1 2 2 2
7 8
1 1 4 3 4 1 3 4
5 10
1 1 1 1 1 1 1 1 1 1
\end{verbatim}

\paragraph{standard output}
\begin{verbatim}
Case #1:
5 4 2
2 3
5 1
2 1 4
3 5 2
1 3 4
Case #2:
6 7
1 3
6 5 1 2 3
7 1 4 2
1 6 4 7 5
7 3
2 6 3 5
3 4 5 2 7
Case #3:
5 3
5 2
4 3
1 5
1 3
2 3
4 2
4 1
1 2
4 5
\end{verbatim}

\section{Problem E. Median}

\paragraph{Tệp vào:} standard input
\paragraph{Tệp ra:} standard output
\paragraph{Giới hạn thời gian:} 1 second
\paragraph{Giới hạn bộ nhớ:} 512 megabytes

Ông docriz có \(n\) số nguyên khác nhau \(1,2,\ldots,n\). Ông muốn chia các
số này thành \(m\) tập rời nhau sao cho trung vị của tập thứ \(j\) là \(b_j\).
Hãy giúp ông xác định liệu điều đó có thể thực hiện được hay không.

Chú ý: Với một tập có kích thước \(k\), sắp xếp các phần tử của nó thành
\(c_1,c_2,\ldots,c_k\). Trung vị của tập này được định nghĩa là
\[
  c_{\lfloor (k+1)/2\rfloor}.
\]

\subsection*{Dữ liệu vào}

Dòng đầu tiên chứa một số nguyên \(T\) (\(1\le T\le 1000\)) là số bộ dữ liệu.
Sau đó là \(T\) bộ dữ liệu.

Dòng đầu tiên của mỗi bộ dữ liệu chứa hai số nguyên \(n,m\)
(\(1\le m\le n\le 10^5\)) là số lượng số nguyên mà ông docriz có và số tập
ông muốn chia thành.

Dòng tiếp theo chứa \(m\) số nguyên \(b_1,b_2,\ldots,b_m\) (\(1\le b_i\le
n\)). Bảo đảm rằng mọi số trong \(b\) đôi một phân biệt.

Bảo đảm rằng \(\sum n\le 2\times 10^6\).

\subsection*{Dữ liệu ra}

Với mỗi bộ dữ liệu, in \texttt{YES} nếu có thể đạt được mục tiêu của ông,
hoặc \texttt{NO} nếu không.

\subsection*{Ví dụ}

\paragraph{standard input}
\begin{verbatim}
3
4 4
2 4 3 1
4 3
1 3 4
4 3
2 3 4
\end{verbatim}

\paragraph{standard output}
\begin{verbatim}
YES
YES
NO
\end{verbatim}

\section{Problem F. The Struggle}

\paragraph{Tệp vào:} standard input
\paragraph{Tệp ra:} standard output
\paragraph{Giới hạn thời gian:} 4 seconds
\paragraph{Giới hạn bộ nhớ:} 512 megabytes

nocriz là một học sinh, và cũng như nhiều người khác, cậu có một ước mơ cho
cuộc đời mình. Không may là cuộc sống không phải lúc nào cũng dễ dàng, và có
những lúc ước mơ dường như rất xa vời còn hành trình theo đuổi thì thật khó
khăn.

Thực tế thường khiến con người rời xa việc theo đuổi ước mơ để chấp nhận
những gì đang có trước mắt. Bị cám dỗ từ bỏ cuộc đấu tranh, nocriz nói với
Artemisia: ``Nhưng tớ không nên từ bỏ ước mơ, đúng không?''. Artemisia trả
lời: ``Tất nhiên là không nên dễ dàng từ bỏ ước mơ! Hãy đấu tranh, đấu tranh
cho đến khi bị tảng đá nghiền nát!''.

Đêm đó, nocriz có một giấc mơ rất tệ. Trong mơ có một tảng đá hình elip, và
cậu được yêu cầu tính giá trị
\[
  \sum (x\oplus y)^3 x^{-2}y^{-1}\bmod 10^9+7
\]
trên tất cả các điểm nguyên \((x,y)\) nằm trong elip, trong đó \(\oplus\) là
phép XOR bit.

Nói hình thức hơn, cho các số nguyên \(a,b,c,d,e,f\), bạn cần tính
\[
  \sum_{(x,y)\in E}(x\oplus y)^3 x^{-2}y^{-1}\bmod 10^9+7,
\]
trong đó
\[
  E=\{(x,y)\mid x,y\in\mathbb{Z},\
  a(x-b)^2+c(y-d)^2+e(x-b)(y-d)\le f\}.
\]
Bảo đảm rằng mọi điểm như vậy thỏa mãn \(0<x,y<4\times 10^6\), và elip chứa
ít nhất một điểm nguyên.

nocriz đã không bị tảng đá nghiền nát ngày hôm đó. Bạn có thể giải bài toán
này như nocriz đã làm không?

\subsection*{Dữ liệu vào}

Dòng đầu tiên chứa một số nguyên \(T\) (\(1\le T\le 10000\)) là số bộ dữ
liệu. Sau đó là \(T\) bộ dữ liệu.

Dòng duy nhất của mỗi bộ dữ liệu chứa sáu số nguyên \(a,b,c,d,e,f\)
\[
  1\le a,c\le 100,\qquad -100\le e\le 100,\qquad
  1\le b,d\le 4\times 10^6,\qquad 0\le f\le 10^{15}.
\]

Bảo đảm rằng có không quá \(1000\) bộ dữ liệu mà
\(\max_{(x,y)\in E}\max(x,y)>100\); không quá \(100\) bộ dữ liệu mà
\(\max_{(x,y)\in E}\max(x,y)>1000\); không quá \(10\) bộ dữ liệu mà
\(\max_{(x,y)\in E}\max(x,y)>10000\); không quá \(3\) bộ dữ liệu mà
\(\max_{(x,y)\in E}\max(x,y)>1000000\); không quá \(1\) bộ dữ liệu mà
\(\max_{(x,y)\in E}\max(x,y)>2000000\), và
\(\max_{(x,y)\in E}\max(x,y)\le 1.1\times 10^7\).

\subsection*{Dữ liệu ra}

Với mỗi bộ dữ liệu, in một số nguyên biểu thị đáp án.

\subsection*{Ví dụ}

\paragraph{standard input}
\begin{verbatim}
2
2 2 1 3 -1 6
13 19 11 17 6 1919
\end{verbatim}

\paragraph{standard output}
\begin{verbatim}
483333436
102810867
\end{verbatim}

\section{Problem G. Power Station of Art}

\paragraph{Tệp vào:} standard input
\paragraph{Tệp ra:} standard output
\paragraph{Giới hạn thời gian:} 3 seconds
\paragraph{Giới hạn bộ nhớ:} 512 megabytes

Là một người yêu nghệ thuật hiện đại, nocriz thích đến Power Station of Art.

Hiện tại, tại Power Station of Art đang trưng bày các tác phẩm thuộc một thể
loại nghệ thuật hiện đại gọi là ``đồ thị đỏ-đen''. Mỗi tác phẩm đồ thị đỏ-đen
là một đồ thị vô hướng có nhãn, với một số và một màu gắn với mỗi đỉnh. Mỗi
đỉnh hoặc màu đỏ hoặc màu đen.

Có thể sửa đổi một đồ thị theo cách sau: chọn một cạnh và hoán đổi các số
được viết trên hai đỉnh tương ứng. Ngoài ra, nếu hai đỉnh có cùng màu, màu
của cả hai đỉnh sẽ bị đổi (từ đỏ sang đen hoặc từ đen sang đỏ). Nếu không,
màu của hai đỉnh được giữ nguyên.

Bây giờ, nocriz đang nghiên cứu hai tác phẩm. Hai đồ thị giống nhau, nhưng
các số và màu có thể khác nhau. Có thể thực hiện một số (có thể bằng không)
lần sửa đổi trên các tác phẩm để khiến chúng trở nên giống nhau hay không?

\subsection*{Dữ liệu vào}

Dòng đầu tiên chứa một số nguyên \(T\) (\(1\le T\le 30000\)) là số bộ dữ
liệu. Sau đó là \(T\) bộ dữ liệu.

Dòng đầu tiên của mỗi bộ dữ liệu chứa hai số nguyên \(n,m\)
(\(1\le n,m\le 10^6\)) là số đỉnh và số cạnh.

Sau đó là \(m\) dòng, mỗi dòng chứa hai số nguyên \(u_i,v_i\)
(\(1\le u_i,v_i\le n,\ u_i\ne v_i\)) biểu thị một cạnh. Bảo đảm rằng dữ liệu
vào không có đa cạnh, và đồ thị có thể không liên thông.

Sau đó là các số và màu của hai đồ thị. Với mỗi đồ thị:

Dòng đầu tiên chứa \(n\) số nguyên, số nguyên thứ \(i\) là \(a_i\)
(\(0\le a_i\le 10^6\)), biểu thị số của đỉnh thứ \(i\) trong đồ thị.

Dòng thứ hai chứa \(n\) ký tự. Nếu ký tự thứ \(i\) là \texttt{R}, đỉnh thứ
\(i\) màu đỏ. Nếu ký tự thứ \(i\) là \texttt{B}, đỉnh thứ \(i\) màu đen.

Bảo đảm rằng \(\sum n\le 10^6\) và \(\sum m\le 10^6\).

\subsection*{Dữ liệu ra}

Với mỗi bộ dữ liệu, in \texttt{YES} nếu có thể thực hiện một số (có thể bằng
không) lần sửa đổi trên các tác phẩm để khiến chúng giống nhau, hoặc
\texttt{NO} nếu không.

\subsection*{Ví dụ}

\paragraph{standard input}
\begin{verbatim}
3
2 1
1 2
3 4
RR
4 3
BB
3 2
1 2
2 3
1 1 1
RBR
1 1 1
BBB
3 3
1 2
2 3
3 1
1 1 1
RBR
1 1 1
BBB
\end{verbatim}

\paragraph{standard output}
\begin{verbatim}
YES
NO
YES
\end{verbatim}

\section{Problem H. Command and Conquer: Red Alert 2}

\paragraph{Tệp vào:} standard input
\paragraph{Tệp ra:} standard output
\paragraph{Giới hạn thời gian:} 8 seconds
\paragraph{Giới hạn bộ nhớ:} 512 megabytes

Là một chàng trai hoài cổ, nocriz thích xem HBK08 và Lantian28 chơi trò
Command and Conquer: Red Alert 2. Tuy nhiên, bản thân cậu lại không biết chơi
trò này.

Trong trò chơi, ban đầu bạn có một lính bắn tỉa ở vị trí
\((-10^{100},-10^{100},-10^{100})\) trong thế giới 3D, và có \(n\) lính địch,
trong đó lính thứ \(i\) ở vị trí \((x_i,y_i,z_i)\). Ta nói tầm bắn của lính
bắn tỉa là \(k\) nếu lính bắn tỉa có thể tiêu diệt mọi kẻ địch thỏa mãn
\[
  \max(|x_s-x_e|,|y_s-y_e|,|z_s-z_e|)\le k,
\]
trong đó \((x_s,y_s,z_s)\) là vị trí của lính bắn tỉa và
\((x_e,y_e,z_e)\) là vị trí của kẻ địch.

Nếu lính bắn tỉa chỉ có thể di chuyển từ \((x,y,z)\) đến \((x+1,y,z)\),
\((x,y+1,z)\) hoặc \((x,y,z+1)\), giá trị nhỏ nhất của \(k\) để lính bắn tỉa
có thể tiêu diệt mọi kẻ địch là bao nhiêu?

Lính bắn tỉa được phép di chuyển số bước không giới hạn và được phép tiêu
diệt kẻ địch bất cứ khi nào anh ta đang ở một tọa độ nguyên.

\subsection*{Dữ liệu vào}

Dòng đầu tiên chứa một số nguyên \(T\) (\(1\le T\le 50000\)) là số bộ dữ
liệu. Sau đó là \(T\) bộ dữ liệu.

Dòng đầu tiên của mỗi bộ dữ liệu chứa một số nguyên \(n\) (\(1\le n\le
5\times 10^5\)) là số kẻ địch.

Sau đó là \(n\) dòng, mỗi dòng chứa ba số nguyên \(x_i,y_i,z_i\)
(\(-10^9\le x_i,y_i,z_i\le 10^9\)) là vị trí của kẻ địch thứ \(i\).

Bảo đảm rằng \(\sum n\le 2\times 10^6\).

\subsection*{Dữ liệu ra}

Với mỗi bộ dữ liệu, in một số nguyên biểu thị đáp án.

\subsection*{Ví dụ}

\paragraph{standard input}
\begin{verbatim}
3
2
0 0 0
1 1 1
2
0 1 0
1 0 1
5
1 1 4
5 1 4
1 9 1
9 8 1
0 0 0
\end{verbatim}

\paragraph{standard output}
\begin{verbatim}
0
1
2
\end{verbatim}

\section{Problem I. Typing Contest}

\paragraph{Tệp vào:} standard input
\paragraph{Tệp ra:} standard output
\paragraph{Giới hạn thời gian:} 1 second
\paragraph{Giới hạn bộ nhớ:} 512 megabytes

Thầy docriz đang lên kế hoạch chọn một số học sinh trong lớp cho một cuộc thi
đánh máy.

Có \(n\) học sinh trong lớp. Tốc độ đánh máy ban đầu của bạn thứ \(i\) là
\(s_i\), và tiếng ồn khi đánh máy là \(f_i\). Tuy nhiên, khi chọn nhiều học
sinh cùng thi, tổng tốc độ đánh máy của họ không phải là tổng tốc độ ban đầu
của từng người, vì tiếng ồn do mỗi người tạo ra ảnh hưởng đến những người
khác.

Cụ thể, nếu các học sinh \(1,2,3,\ldots,k\) tạo thành một đội, tốc độ đánh
máy thực tế của học sinh \(1\) là
\[
  s_1(1-f_1f_2-f_1f_3-\cdots-f_1f_k),
\]
tốc độ đánh máy thực tế của học sinh \(2\) là
\[
  s_2(1-f_2f_1-f_2f_3-\cdots-f_2f_k),
\]
và tương tự.

Thầy docriz muốn lập một đội sao cho tổng tốc độ đánh máy lớn nhất có thể.
Hãy giúp thầy tính tốc độ đánh máy tối đa có thể đạt được.

\subsection*{Dữ liệu vào}

Dòng đầu tiên chứa một số nguyên \(T\) (\(1\le T\le 1000\)) là số bộ dữ
liệu. Sau đó là \(T\) bộ dữ liệu.

Dòng đầu tiên của mỗi bộ dữ liệu chứa một số nguyên \(n\) (\(1\le n\le 100\))
là số học sinh.

Sau đó là \(n\) dòng, mỗi dòng chứa hai số \(s_i,f_i\)
(\(1\le s_i\le 10^{12},\ 0\le f_i\le 1\)), trong đó \(s_i\) là số nguyên và
\(f_i\) là số thực có đúng \(2\) chữ số sau dấu thập phân.

Bảo đảm rằng
\[
  \sum n<1.2\times 10^4,\qquad
  \sum n^2<2.4\times 10^5,\qquad
  \sum n^3<1.2\times 10^7,\qquad
  \sum n^4<9\times 10^8.
\]

\subsection*{Dữ liệu ra}

Với mỗi bộ dữ liệu, in một số thực là tốc độ đánh máy tối đa mà thầy docriz
có thể đạt được. Hãy giữ đúng \(9\) chữ số sau dấu thập phân.

Bảo đảm rằng đáp án không có sai số độ chính xác khi giữ \(9\) chữ số sau dấu
thập phân, nên bài này không dùng trình chấm đặc biệt. Hãy bảo đảm độ chính
xác của kết quả in ra.

\subsection*{Ví dụ}

\paragraph{standard input}
\begin{verbatim}
4
3
10 0.00
11 0.00
12 0.00
3
10 1.00
11 1.00
12 1.00
3
10 0.50
11 0.50
12 0.50
3
10 0.33
11 0.21
12 0.92
\end{verbatim}

\paragraph{standard output}
\begin{verbatim}
33.000000000
12.000000000
17.250000000
20.421900000
\end{verbatim}

\subsection*{Ghi chú}

Nhắc nhở thân thiện: trình chấm chạy rất nhanh. Khi đánh giá liệu thuật toán
bạn thiết kế có thể vượt qua hay không, nên tham khảo thông tin sau.

Tích chập đa thức với độ phức tạp thời gian \(O(n^2)\) mất khoảng \(1\) giây
trên trình chấm này với \(n=40000\).

Mã kiểm thử có tại đây: \texttt{https://www.luogu.com.cn/paste/zq0kp7pm}.

\section{Problem J. Array}

\paragraph{Tệp vào:} standard input
\paragraph{Tệp ra:} standard output
\paragraph{Giới hạn thời gian:} 3 seconds
\paragraph{Giới hạn bộ nhớ:} 512 megabytes

Koishi cho bạn một mảng \(B\) độ dài \(n\) thỏa mãn
\[
  1\le B_1\le B_2\le \cdots\le B_n\le n+1.
\]
Gọi \(S(T)\) là tập các số xuất hiện trong mảng \(T\). Cô ấy hỏi liệu có tồn
tại một mảng \(A\) độ dài \(n\) sao cho với mọi \(l,r\) (\(1\le l\le r\le
n\)), \(S(A[l,r])=S(A[1,n])\) đúng khi và chỉ khi \(r\ge B_l\) hay không.

\(A[l,r]\) biểu thị mảng con của \(A\) tạo bởi \(A_l,A_{l+1},\ldots,A_r\).

Chú ý: Nếu tồn tại \(i\) (\(1\le i\le n\)) sao cho \(B_i<i\), thì mảng \(A\)
cần tìm chắc chắn không tồn tại.

\subsection*{Dữ liệu vào}

Dòng đầu tiên chứa một số nguyên \(T\) (\(1\le T\le 3\times 10^4\)) là số bộ
dữ liệu. Sau đó là \(T\) bộ dữ liệu.

Dòng đầu tiên của mỗi bộ dữ liệu chứa một số nguyên \(n\) (\(1\le n\le
2\times 10^5\)) là độ dài của mảng \(B\) và \(A\).

Dòng tiếp theo chứa \(n\) số nguyên \(B_1,B_2,\ldots,B_n\)
\[
  1\le B_1\le B_2\le \cdots\le B_n\le n+1,
\]
là mảng mà Koishi cho bạn.

Bảo đảm rằng \(\sum n\le 4.5\times 10^6\).

\subsection*{Dữ liệu ra}

Với mỗi bộ dữ liệu, in \texttt{YES} nếu mảng \(A\) tồn tại, hoặc \texttt{NO}
nếu không.

\subsection*{Ví dụ}

\paragraph{standard input}
\begin{verbatim}
3
4
3 3 5 5
7
4 6 6 7 8 8 8
5
2 3 4 4 6
\end{verbatim}

\paragraph{standard output}
\begin{verbatim}
YES
YES
NO
\end{verbatim}

\section{Problem K. Game}

\paragraph{Tệp vào:} standard input
\paragraph{Tệp ra:} standard output
\paragraph{Giới hạn thời gian:} 6 seconds
\paragraph{Giới hạn bộ nhớ:} 512 megabytes

Koishi đang chơi một trò chơi với Satori.

Có một mảng độ dài \(10^{18}\). Trong trò chơi, Koishi và Satori luân phiên
thao tác trên mảng này, và Koishi đi trước. Khi đến lượt một người chơi, nếu
trong mảng chỉ còn một phần tử, cô ấy thua ngay lập tức. Nếu không, cô ấy cần
xóa số ngoài cùng bên trái hoặc số ngoài cùng bên phải của mảng còn lại.

Trò chơi quá nhàm chán với Koishi, nên cô ấy nghĩ ra các luật sau.

Có \(n\) đoạn con của mảng này là đặc biệt. Cụ thể, đoạn con thứ \(i\) được
mô tả bởi ba số nguyên \((l_i,r_i,z_i)\), nghĩa là khi đến lượt một người
chơi, nếu mảng còn lại là đoạn con \([l_i,r_i]\), cô ấy sẽ thắng ngay nếu
\(z_i=1\) hoặc thua ngay nếu \(z_i=0\). Hơn nữa, nếu có một đoạn con đặc biệt
\((x,x,1)\), thì người chơi sẽ thắng ngay lập tức khi mảng còn lại là
\([x,x]\) (trong trường hợp mặc định, khi chỉ còn một số, người chơi sẽ
thua).

Sẽ có \(q\) ván chơi. Khi bắt đầu ván thứ \(i\), Utuoho sẽ đưa cho hai người
chơi đoạn con \([a_i,b_i]\) và lấy đi tất cả các phần khác của mảng. Nghĩa
là Koishi và Satori chỉ chơi trên đoạn con \([a_i,b_i]\). Tất cả \(q\) ván
chơi độc lập với nhau.

Hai người chơi luôn sử dụng chiến lược tối ưu. Hãy cho họ biết ai sẽ thắng
trong mỗi ván.

\subsection*{Dữ liệu vào}

Dòng đầu tiên chứa một số nguyên \(T\) (\(1\le T\le 2\times 10^3\)) là số bộ
dữ liệu. Sau đó là \(T\) bộ dữ liệu.

Dòng đầu tiên của mỗi bộ dữ liệu chứa hai số nguyên \(n,q\)
(\(1\le n,q\le 10^5\)) là số đoạn con đặc biệt và số ván chơi.

Sau đó là \(n\) dòng, mỗi dòng chứa ba số nguyên \(l_i,r_i,z_i\)
\[
  1\le l_i\le r_i\le 10^9,\qquad 0\le z_i\le 1,
\]
thỏa mãn rằng với mọi \(i\ne j\), \((l_i,r_i)\ne(l_j,r_j)\).

Sau đó là \(q\) dòng, mỗi dòng chứa hai số nguyên \(a_i,b_i\)
(\(1\le a_i\le b_i\le 10^9\)) là đoạn con ban đầu của ván thứ \(i\).

Bảo đảm rằng \(\sum (n+q)\le 9\times 10^5\).

\subsection*{Dữ liệu ra}

Với mỗi bộ dữ liệu, in một dòng gồm \(q\) số nguyên \(v_i\) (\(0\le v_i\le
1\)) không có dấu cách, trong đó \(v_i=0\) nếu Koishi thua ở ván thứ \(i\) và
\(v_i=1\) nếu cô ấy thắng.

\subsection*{Ví dụ}

\paragraph{standard input}
\begin{verbatim}
1
5 10
1 2 0
3 3 1
3 4 1
2 4 1
1 3 0
1 2
1 3
1 4
1 5
2 3
2 4
2 5
3 4
3 5
4 5
\end{verbatim}

\paragraph{standard output}
\begin{verbatim}
0010111101
\end{verbatim}

\end{document}
